\documentclass[11pt]{article}
\usepackage[margin=1in]{geometry}
\usepackage{amsmath,amssymb,mathtools}
\usepackage[T1]{fontenc}
\usepackage{lmodern}
\usepackage{enumitem}
\usepackage{microtype}

\title{Mathematical Model of the Implemented Multi-AGV Scheduling Problem}
\author{}
\date{}

\begin{document}
\maketitle

\section{Sets and graph}

Let $G=(N,E)$ be the undirected factory graph. Each undirected edge represents a
bidirectional track. Let $F$ be the set of floors, and let $N_f\subseteq N$
denote the nodes on floor $f\in F$. Let $L\subseteq E$ be the set of lift links
between floors. The parking node is $\pi\in N$.

Let $J$ be the set of jobs and $V$ the set of AGVs. Each job $j\in J$ is served
by exactly one AGV. Let $M_i$ index the successive node visits on AGV $i$'s
route, including repeated visits to the same physical node.

\section{Parameters}

For every job $j\in J$, its pickup and delivery nodes are
\[
  p_j\in N,\qquad p'_j\in N,\qquad p_j\ne p'_j,
\]
and $t_j,t'_j$ are its desired pickup and delivery times. The service-time
parameters are $\eta$ (loading) and $\eta'$ (unloading). The safety separation
is $\delta>0$, and AGV speed is $v>0$. Each edge $e\in E$ has length $\ell_e$;
for lift edges $e\in L$, $\ell_e$ is the configured lift length. Travel time
along an edge is
\[
  \tau_{uv}=\frac{\ell_{\{u,v\}}}{v},\qquad \{u,v\}\in E.
\]
Let $d(u,w)$ be the shortest-path travel time between nodes $u,w$, computed
from these edge times. The code uses a fixed weight $\lambda\ge 0$ in its
assignment fitness and a maximum of $k_{\max}$ passes in Algorithm 3.

Jobs are considered in the release order
\[
  j_1,\ldots,j_{|J|}
  =\operatorname{sort}_{j\in J}(t_j,j),
\]
that is, increasing desired pickup time, with job identifier as the tie-breaker.
AGV $i$ has a fixed priority rank $\rho_i$, where smaller values mean higher
priority.

\section{Decision and schedule variables}

The assignment variable is
\[
  x_{ji}\in\{0,1\},
\]
where $x_{ji}=1$ means AGV $i$ is assigned job $j$. Each job is assigned once:
\[
  \sum_{i\in V}x_{ji}=1\qquad \forall j\in J.
\]

For each AGV $i$, let
\[
  R_i=(n_{i,0},n_{i,1},\ldots,n_{i,q_i})
\]
be its ordered route, with $n_{i,0}=n_{i,q_i}=\pi$. The implemented planner
constructs $R_i$ by concatenating shortest paths from the AGV's current node
to each assigned job's pickup and delivery nodes, then back to $\pi$.

For visit $m\in M_i$, let $lt_{im}$ and $rt_{im}$ be the start and end of its
node-occupancy window. Its dwell time is $dwell_{im}$, equal to $\eta$ at a
pickup, $\eta'$ at a delivery, and zero at a pass-through visit. The edge
occupancy window for consecutive visits $m-1,m$ is
\[
  [rt_{i,m-1},lt_{im}],
\]
with duration fixed to the edge travel time:
\[
  lt_{im}-rt_{i,m-1}=\tau_{n_{i,m-1},n_{i,m}}.
\]
Thus, the AGV does not stop while traversing an edge. At ordinary nodes,
\[
  rt_{im}-lt_{im}=dwell_{im}.
\]
The parking node is exempt from node-conflict checks and may have an extended
window to represent waiting.

\section{Algorithm 1 assignment quantities}

When job $j$ is considered, let $s_i$ and $vt_i$ be AGV $i$'s current node and
available time. The code computes
\begin{align*}
  ct_{ji}
  &=vt_i+d(s_i,p_j)+\eta+d(p_j,p'_j)+\eta',\\
  et_{ji}
  &=\max\!\left(0,\ t_j-\bigl(vt_i+d(s_i,p_j)+\eta\bigr)\right),\\
  dt_{ji}
  &=\max(0,\ ct_{ji}-t'_j),\\
  F_{ji}
  &=ct_{ji}+\lambda(et_{ji}+dt_{ji}).
\end{align*}
The job is assigned to the AGV minimizing $F_{ji}$; ties are broken by AGV
priority. After assignment, $vt_i$ and $s_i$ are updated to $ct_{ji}$ and
$p'_j$, respectively.

\section{Conflict separation and start delays}

Let $W_n$ be the fixed occupancy windows of higher-priority AGVs at ordinary
node $n$. Let $W'_{uv}$ be the fixed edge windows on directed traversal
$u\to v$. In the implementation, a physical track is exclusive in both
directions: edge conflicts are checked against windows for both $u\to v$ and
$v\to u$.

For two different AGVs $i\ne h$, overlapping use of an ordinary node or the
same bidirectional track must be separated by at least $\delta$. For occupancy
intervals $[a,b]$ and $[c,d]$, the required separation is
\[
  c\ge b+\delta
  \quad\text{or}\quad
  a\ge d+\delta.
\]
The parking node $\pi$ is exempt from node separation. Algorithm 3 resolves a
conflict by shifting a prefix of the current route to the right. It computes
the cumulative start delay $st^*$ by sweeping backward from the conflict visit
to the nearest previous exempt node, using the window function
$f(lt,rt,W)$. After the sweep, the affected node and edge windows are shifted
by $st^*$, and the previous parking window is extended by the same amount. The
sweep is limited to $k_{\max}$ passes; if exhausted, the implementation applies
its fallback delay.

\section{Interpretation and assumptions}

This is the scheduling procedure implemented by \texttt{agv\_core.py} over the
layouts in \texttt{agv\_layouts.py}, rather than a single globally optimized
mixed-integer program. It first assigns jobs greedily in pickup-release order
using $F_{ji}$, fixes shortest-path routes, then schedules AGVs in priority
order and resolves conflicts through Algorithm 3. It assumes known job times,
fixed AGV priorities, constant speed, fixed edge travel times, service time
only at pickup and delivery, no stops on edges, and exclusive use of each
bidirectional track. The implemented centralized model allows waiting only at
$\pi$; lift-lobby waiting points are an additional hybrid-layout feature.

\section{Hard feasibility constraints}

The following conditions express the hard checks performed by the independent
\texttt{verify()} routine. Let $P_{im}=[lt_{im},rt_{im}]$ be the window for
visit $m$ of AGV $i$, and let $\mathcal{B}$ be the set of exempt waiting nodes.
For the centralized layout, $\mathcal{B}=\{\pi\}$; in the hybrid layout it
also contains lift lobbies. Write $\varepsilon$ for the verifier's numerical
tolerance.

\paragraph{Route endpoints and initial time.}
Every AGV starts at the parking node at time zero and finishes at that node:
\[
  n_{i,0}=\pi,\qquad n_{i,q_i}=\pi,\qquad |lt_{i,0}|\le\varepsilon
  \qquad (i\in V).
\]

\paragraph{Node conflict avoidance.}
For each ordinary node $n\in N\setminus\mathcal{B}$, distinct AGVs may not
occupy it within the safety separation. For every pair of visits $(i,m)$ and
$(h,r)$ to $n$, where $i\ne h$, the windows must satisfy
\[
  rt_{im}+\delta\le lt_{hr}
  \quad\lor\quad
  rt_{hr}+\delta\le lt_{im}.
\]
No node-conflict constraint is imposed on $n\in\mathcal{B}$: parking and
lobbies are exempt and may hold multiple waiting AGVs.

\paragraph{Track conflict avoidance in both directions.}
For each physical track $e=\{u,v\}\in E$, define the traversal window for a
visit transition $m-1\to m$ by
\[
  Q_{im}=[a_{im},b_{im}]
  =[rt_{i,m-1},lt_{im}],
  \qquad \{n_{i,m-1},n_{i,m}\}=e.
\]
For every pair of traversals of the same physical track by distinct AGVs,
including traversals in opposite directions, require
\[
  b_{im}+\delta\le a_{hr}
  \quad\lor\quad
  b_{hr}+\delta\le a_{im}.
\]
This is the verifier's bidirectional-track rule: use of $u\to v$ conflicts
with both $u\to v$ and $v\to u$ use by another AGV.

\paragraph{No pause on an edge.}
Every edge window has exactly the edge's non-stop travel duration:
\[
  lt_{im}-rt_{i,m-1}
  =\tau_{n_{i,m-1},n_{i,m}}
  =\frac{\ell_{\{n_{i,m-1},n_{i,m}\}}}{v}
  \qquad (m=1,\ldots,q_i).
\]

\paragraph{Dwell times.}
Let $\mathcal{P}_{im}$ and $\mathcal{D}_{im}$ indicate that visit $(i,m)$
performs a pickup or delivery, respectively. With at most one service event
per visit, define
\[
  d_{im}=\eta\,\mathcal{P}_{im}+\eta'\,\mathcal{D}_{im}.
\]
At every non-exempt node, the occupancy window equals the required service
dwell (and is zero at pass-through visits):
\[
  rt_{im}-lt_{im}=d_{im}
  \qquad(n_{i,m}\notin\mathcal{B}).
\]
At an exempt parking or lobby visit, the verifier only requires the window to
be at least as long as any service dwell there:
\[
  rt_{im}-lt_{im}\ge d_{im}
  \qquad(n_{i,m}\in\mathcal{B}),
\]
so the excess duration can represent waiting. In the implemented job routes,
$\eta$ is used only at pickup events, $\eta'$ only at delivery events, and
pass-through visits have zero dwell.

\paragraph{Exactly-once job service.}
Let $z_{ij}$ be the number of times job $j$ appears in AGV $i$'s assigned-job
list. Every job must be assigned to exactly one AGV:
\[
  \sum_{i\in V}z_{ij}=1\qquad(j\in J),\qquad z_{ij}\in\mathbb{Z}_{\ge0}.
\]
Equivalently, the multiset union of all AGV job lists must equal the job set,
with each job occurring once. In the implementation, \texttt{verify()} checks
this condition by comparing the flattened AGV job lists with the job IDs; it
does not independently verify that each listed job's pickup and delivery
events occur at the corresponding stations and in the correct order.

\section{Algorithm 3 start-delay conditions}

For a conflict at visit $p$, Algorithm 3 chooses a nonnegative cumulative
right-shift $st^*\ge0$ for the route segment after the nearest preceding
exempt node $q$ through $p$. In the no-fallback case, its update is
\[
  lt'_{im}=lt_{im}+st^*,\qquad rt'_{im}=rt_{im}+st^*
  \quad(q<m\le p),\qquad
  rt'_{iq}=rt_{iq}+st^*.
\]
The last equality extends the waiting window at the exempt node. Edge windows
within the shifted segment move consistently because both endpoint times are
shifted; the connecting edge from $q$ is delayed by extending $rt_{iq}$. The
algorithm accumulates the smallest right-shifts needed to clear the fixed
higher-priority windows encountered in its backward sweep, while checking
that a feasible shift does not cross the next blocking window (the $mt^*$
bound). It repeats up to $k_{\max}$ passes. If those passes are exhausted,
the code uses its fallback shift so the AGV waits at the previous exempt node
until higher-priority traffic is free. The independent verifier subsequently
checks the resulting complete plan.

\section{Reported objectives}

The implementation reports three performance measures. Let $C_j$ be the
actual delivery-completion time of job $j$ in the final plan, and let $U_i$
be AGV $i$'s usage time, defined in the code as its return time to $\pi$ when
it serves jobs and zero otherwise. Then
\begin{align*}
  M &= \max_{j\in J} C_j
      &&\text{(makespan, last delivery)},\\
  T &= \sum_{j\in J}\max(0,C_j-t'_j)
      &&\text{(total delivery tardiness)},\\
  I &= \max_{i\in V}U_i-\min_{i\in V}U_i
      &&\text{(load imbalance)}.
\end{align*}
One may state a scalar objective as
\[
  \min\; \alpha M+\beta T+\gamma I,
  \qquad \alpha,\beta,\gamma\ge0,
\]
or a lexicographic objective as
\[
  \operatorname{lexmin}(M,T,I).
\]
These are possible mathematical formulations for comparing schedules; the
code does not choose weights or optimize either combined objective. It uses
$F_{ji}$ greedily for job assignment, then applies prioritized conflict
resolution, and reports $M$, $T$, and $I$ as metrics.

\paragraph{Enforced versus checked conditions.}
Algorithm 2 constructs each route's node and edge windows and calls Algorithm
3 when a new window conflicts with already frozen higher-priority windows.
Algorithm 3 shifts the affected route prefix, extends the preceding exempt
node's window, and thereby enforces conflict avoidance incrementally; in the
hybrid planner, floor servers also plan legs and the central server reserves
lift tracks. The independent \texttt{verify()} routine does not plan or repair
the schedule. It checks the completed plans for start/end at $\pi$, node and
bidirectional-track separation, exact edge travel durations, dwell-window
lengths, and exactly-once job service. Thus, the full-plan conditions are
validated by \texttt{verify()} even though the construction relies on
Algorithm 2/3 to satisfy them.

\end{document}
